\subsection{Abstand eines Punktes von einer Geraden}
\label{subsec:analytic_geometry_point_to_line_distance}

Die kürzeste Verbindung eines Punktes $Q$ mit einer Geraden $g$ steht
senkrecht auf der Geraden.

\begin{center}
  \CachedTikzFigure{content/resources/chapters/analytic_geometry/tikz/point_line_distance.tex}
\end{center}

\begin{TheoremBox}{Punkt--Geraden-Abstand}{analytic_geometry_point_line_distance}
  Sei $P$ ein Punkt der Geraden $g$ und sei $Q$ ein beliebiger Punkt. Mit
  einem Einheitsrichtungsvektor $\vec e_a$ von $g$ ist der parallele Anteil
  des Verbindungsvektors
  \[
    \overrightarrow{PQ}_{\parallel}
    =(\overrightarrow{PQ}\cdot\vec e_a)\vec e_a.
  \]
  Daher entspricht der Abstand $d$ dem Betrag der senkrechten Komponente
  $\overrightarrow{PQ}_{\perp}$ von $\overrightarrow{PQ}$:
  \[
    d
    =\left\lvert\overrightarrow{PQ}_{\perp}\right\rvert
    =\left\lvert
      \overrightarrow{PQ}-\overrightarrow{PQ}_{\parallel}
    \right\rvert
    =\left\lvert
      \overrightarrow{PQ}-(\overrightarrow{PQ}\cdot\vec e_a)\vec e_a
    \right\rvert.
  \]
\end{TheoremBox}

\begin{ExampleBox}{Abstand eines Punktes von einer Geraden}{analytic_geometry_point_line_distance_example}
  Gegeben sind die Gerade $g$ in Parameterform und der Punkt $Q(0,2)$:
  \[
    g:\quad\vec x=\overrightarrow{OP}+\lambda\vec a,
    \qquad
    \overrightarrow{OP}=\begin{pmatrix}2\\2\end{pmatrix},
    \quad
    \vec a=\begin{pmatrix}1\\-1\end{pmatrix}.
  \]
  Zunächst bestimmen wir den Verbindungsvektor und den Einheitsrichtungsvektor:
  \begin{ClozeArea}[34]{3}
    \[
      \overrightarrow{PQ}=\overrightarrow{OQ}-\overrightarrow{OP}
        =\begin{pmatrix}0\\2\end{pmatrix}-\begin{pmatrix}2\\2\end{pmatrix}
        =\begin{pmatrix}-2\\0\end{pmatrix},
      \qquad
      \vec e_a=\frac{\vec a}{\lvert\vec a\rvert}
        =\frac1{\sqrt2}\begin{pmatrix}1\\-1\end{pmatrix}.
    \]
  \end{ClozeArea}
  Dann bestimmen wir die parallele und die senkrechte Komponente von $\overrightarrow{PQ}$:
  \begin{itemize}
    \item Paralleler Anteil:
      \begin{ClozeArea}[32]{3}
        \[
          \overrightarrow{PQ}_{\parallel}
            =(\overrightarrow{PQ}\cdot\vec e_a)\vec e_a
            =-\sqrt2\,\vec e_a
            =\begin{pmatrix}-1\\1\end{pmatrix}.
        \]
      \end{ClozeArea}
    \item Senkrechter Anteil:
      \begin{ClozeArea}[32]{3}
        \[
          \overrightarrow{PQ}_{\perp}
            =\overrightarrow{PQ}-\overrightarrow{PQ}_{\parallel}
            =\begin{pmatrix}-2\\0\end{pmatrix}-\begin{pmatrix}-1\\1\end{pmatrix}
            =\begin{pmatrix}-1\\-1\end{pmatrix}.
        \]
      \end{ClozeArea}
  \end{itemize}
  \noindent
  \begin{minipage}[c]{0.44\linewidth}
    Den Abstand $d$ erhalten wir gemäß
    Satz~\ref{thm:analytic_geometry_point_line_distance} als Betrag der
    senkrechten Komponente von $\overrightarrow{PQ}$:
    \begin{ClozeArea}[15]{4}
      \[
        \begin{aligned}
          d&=\left\lvert\overrightarrow{PQ}_{\perp}\right\rvert
            =\left\lvert\begin{pmatrix}-1\\-1\end{pmatrix}\right\rvert\\[0.6em]
           &=\sqrt{(-1)^2+(-1)^2}=\sqrt2.
        \end{aligned}
      \]
    \end{ClozeArea}
  \end{minipage}\hfill
  \begin{minipage}[c]{0.54\linewidth}
    \centering
    \CachedTikzFigure{content/resources/chapters/analytic_geometry/tikz/point_line_distance_example.tex}
  \end{minipage}
\end{ExampleBox}
